\subsection{A.2: an exact cyclotomic arithmetic criterion}

We now isolate an arithmetic instance of the eventual upper bound. This is
a proved equivalence, not a proof of either the general conjecture in A.2
or Iwasawa's $\mu$-conjecture. We use $\mu(L)=0$ to mean vanishing of the
class-group Iwasawa invariant of the cyclotomic $\Z_p$-field $L$; the
vanishing condition is independent of the finite initial layer.

\begin{lemma}[Proved from known arithmetic inputs; a cyclotomic Dedekind domain]
\label{A:cyclotomic-domain}
Fix a prime $p$ and a number field $F$ containing $\mu_p$; if $p=2$,
assume $i\in F$. Let $K=F_{\mathrm{cyc}}$ be its cyclotomic
$\Z_p$-extension, and put $R=\mathcal O_K[1/p]$. Then $R$ is an excellent
Dedekind domain, $\cd_p(K)=1$, and every residue field at a closed point
of $\Spec R$ has $p$-cohomological dimension zero. The same assertions
hold after replacing $K$ by any finite extension $L/K$ and $R$ by its
normalization $R_L=\mathcal O_L[1/p]$.
\end{lemma}
\begin{proof}
Write $K=\bigcup_n F_n$ with $[F_n:F]=p^n$. At a finite prime
$v\nmid p$ of $F$, the tower is unramified. Its Frobenius has infinite
order in $\operatorname{Gal}(K/F)\simeq\Z_p$: under the cyclotomic
character, a zero image would say that the positive integer $Nv>1$ is
a $p$-adic root of unity. Hence the decomposition group is open, so
only finitely many primes lie above $v$. The ramification indices are
one. Lemma~\ref{A:A2-tower-Dedekind} therefore proves that $R$ is
Dedekind. The cyclotomic ramification statements used here are also
\citet[Proposition 11.1.1]{A_NSW}; at primes above $p$ the inertia
subgroups are open and the tower is eventually totally ramified.
A Dedekind domain whose fraction field has characteristic zero is
excellent by \citet[Proposition 15.53.3(4), Tag 07QS]{A_Stacks}.

Every finite extension $L/K$ has the form $F'K$ for a number field
$F'$ containing $F$. Since $K=F\Q_{\mathrm{cyc}}$, this is exactly
the cyclotomic $\Z_p$-extension of $F'$, after changing its finite
initial layer. Repeating the preceding argument proves the assertions
about $R_L$; excellence also makes $R_L$ finite over $R$.

We check the fraction-field dimension, including the passage to arbitrary
coefficients. For every such $L$, a class of $\operatorname{Br}(L)[p]$
descends to a finite layer. Brauer--Hasse--Noether injects the Brauer
group of that number field into the direct sum of its local Brauer
groups. Only finitely many local invariants occur, and restriction to
a finite local extension multiplies an invariant by the local degree.
At every finite place these local degrees in the cyclotomic tower are
unbounded powers of $p$: at $p$ by ramification and elsewhere by the
nonzero Frobenius just checked. Thus one later layer kills every local
invariant simultaneously. There are no real places. Consequently
\[
 \operatorname{Br}(L)[p]=0
 \quad\hbox{for every finite }L/K.
\]
The precise local-global and degree inputs, including the unbounded
local degrees, are \citet[Chapter VII, Theorem 7.1, Proposition 7.2,
and the proof of Lemma 7.3]{A_MilneCFT}.

Since $K$ is an algebraic extension of the totally imaginary number
field $F$, initially $\cd_p(K)\leq2$ by \citet[Proposition 8.3.18, taking all primes in $S$]{A_NSW} and closed-subgroup monotonicity. If $M$ is a finite discrete
$\F_p[\operatorname{Gal}(K^{\mathrm{sep}}/K)]$-module, choose a finite
Galois extension $L/K$ trivializing it. A sum of regular permutation
modules surjects onto $M$. Its second cohomology vanishes by Shapiro,
Kummer theory, and $\mu_p\subset L$, because
$H^2(L,\F_p)\simeq\operatorname{Br}(L)[p]=0$.
The kernel has no third cohomology by the preliminary bound two.
The long exact sequence gives $H^2(K,M)=0$. Filtered colimits and
coefficient d\'evissage give $\cd_p(K)\leq1$. To see equality, choose
a prime above $v\nmid p$ and an element of $F$ of valuation one at $v$.
Its valuation remains one in $K$, so it is not a $p$th power. Its Kummer
class is nonzero in $H^1(K,\mu_p)$.

Finally, a residue field of $R$ is the union of finite fields whose
successive degrees are eventually $p$. Its absolute Galois group has
no nontrivial pro-$p$ subgroup, so its $p$-cohomological dimension is
zero. This also follows from the valuation inequality
$1+\cd_p\kappa(v)\leq\cd_p(K)=1$ already proved in Part A.
\end{proof}

\begin{lemma}[Proved from Iwasawa's structure theorem; the persistent obstruction]
\label{A:cyclotomic-Pic}
Use the notation of Lemma~\ref{A:cyclotomic-domain}. For every nonempty
open $U\subset\Spec R$,
\[
 H^2(U,\mu_p)\simeq\operatorname{Pic}(U)/p.
\]
If $\mu(K)=0$, these groups are zero. If $\mu(K)>0$, they have infinite
dimension over $\F_p$, even after deleting any finite set of closed
points. The assertions hold for every finite extension $L/K$ as well.
\end{lemma}
\begin{proof}
Let $C_K=\varinjlim_n\operatorname{Cl}(\mathcal O_{F_n})$, where the
maps are extension of ideals, and let $A_K$ be its $p$-primary part.
The Iwasawa class-group structure theorem gives the following two
facts about $M=A_K^\vee$, its Pontryagin dual: $M$ is a finitely
generated torsion $\Lambda=\Z_p[[T]]$-module with the class-group
$\mu$-invariant, and it has no nonzero finite $\Lambda$-submodule.
We use \citet[Proposition 3.2.15 and Corollary 3.2.16]{A_SharifiIwasawa},
whose proofs were opened, together with the structure theorem in
Theorem 2.2.13 of that source.

If $\mu(K)=0$, the structure theorem makes $M$ finitely generated
over $\Z_p$. Its $p$-power torsion is then finite, hence zero by the
absence of finite $\Lambda$-submodules. Thus $M$ is free over $\Z_p$
and $A_K$ is $p$-divisible. If $\mu(K)>0$, the elementary module of
$M$ has a summand $\Lambda/(p^a)$ with $a\geq1$. Its subgroup
annihilated by $p$ is infinite. A pseudo-isomorphism changes this
subgroup only by finite kernels and cokernels, so $M[p]$ is infinite.
Duality identifies $M[p]=(A_K/pA_K)^\vee$, proving that $A_K/pA_K$
has infinite $\F_p$-dimension. The prime-to-$p$ part of $C_K$ is
$p$-divisible, hence the same alternatives hold for $C_K/pC_K$.

Invertible modules descend through a filtered union of rings, as do their isomorphisms. Hence the Picard group of $R=\varinjlim_n\mathcal O_{F_n}[1/p]$ is the direct limit of the finite-layer Picard groups. Localization of ideals therefore gives
\[
 \operatorname{Pic}(R)=C_K/D_p,
\]
where $D_p$ is the image of the direct limit of the groups generated
by primes above $p$. After a finite layer these primes no longer
split and are totally ramified at each step. That direct limit is a
finite direct sum of copies of $\Z[1/p]$, so $D_p=pD_p$.
In particular $D_p\subset pC_K$, and
\[
 \operatorname{Pic}(R)/p\simeq C_K/pC_K.
\]
Every nonempty open in the noetherian one-dimensional scheme
$\Spec R$ has finite complement. Its Picard group is the quotient
of $\operatorname{Pic}(R)$ by the classes of those finitely many
primes. Such a quotient preserves vanishing modulo $p$, and it
cannot kill an infinite-dimensional quotient modulo $p$.

For clarity, the Brauer comparison used in Kummer theory also applies
to these possibly non-finite-type Dedekind schemes. For the generic
inclusion $j:\Spec K\to U$, the divisor sequence of étale sheaves is
\[
 0\longrightarrow\Gm\longrightarrow j_*\Gm
   \longrightarrow\bigoplus_{x\in U^{(1)}}i_{x*}\Z\longrightarrow0.
\]
At strictly henselian localizations its exactness is the valuation
sequence of a discrete valuation ring. Hilbert 90 gives
$R^1j_*\Gm=0$, while $H^1(\kappa(x),\Z)=0$. The first cohomology of the displayed direct sum is zero as well: express it as the filtered union of finite sums and use cohomological continuity on the quasi-compact, quasi-separated scheme $U$. The long exact sequence
and the Leray sequence therefore inject
$H^2(U,\Gm)$ into $H^2(K,\Gm)=\operatorname{Br}(K)$.
The latter has no $p$-torsion by Lemma~\ref{A:cyclotomic-domain}.
Kummer theory now gives the asserted isomorphism. The proof for
$L/K$ is identical.
\end{proof}

\begin{theorem}[Proved; A.2 on cyclotomic Dedekind domains is equivalent to a precise arithmetic condition]
\label{A:Iwasawa-equivalence}
For $R$ and $K$ in Lemma~\ref{A:cyclotomic-domain}, the following
conditions are equivalent:
\begin{enumerate}[label=(\roman*)]
\item Some nonempty affine open $U\subset\Spec R$ has $\cd_p(U)\leq1$.
\item Every nonempty affine open $U\subset\Spec R$ has $\cd_p(U)=1$.
\item For every finite extension $L/K$, the cyclotomic class-group
      invariant satisfies $\mu(L)=0$.
\end{enumerate}
Consequently A.2 for this excellent one-dimensional domain holds
if and only if (iii) holds. If (iii) fails, every nonempty affine open
of $\Spec R$ has $p$-cohomological dimension exactly two.
\end{theorem}
\begin{proof}
Let $U$ be any nonempty affine open. Corollary~\ref{A:generic-plus-dimension}
gives the preliminary bound $\cd_p(U)\leq2$. Its closed residue fields
have $p$-cohomological dimension zero.

Suppose (i), and let $f:Y\to U$ be the finite normalization in
a finite extension $L/K$. It is an open arithmetic Dedekind scheme
of the kind in Lemma~\ref{A:cyclotomic-Pic}. Finite pushforward of
torsion étale sheaves is exact and computes cohomology:
\[
 H^q(Y,E)=H^q(U,f_*E).
\]
Thus $\cd_p(Y)\leq\cd_p(U)\leq1$.
If $\mu(L)>0$, Lemma~\ref{A:cyclotomic-Pic} contradicts this with
$E=\mu_p$ and $q=2$. This proves (iii).

Conversely assume (iii). Let $F_0$ be a constructible sheaf of
$\F_p$-vector spaces on $U$. Choose a finite Galois extension $L/K$
trivializing its generic fibre, let $f:Y\to U$ be the normalization,
and choose a dense open $j:V\hookrightarrow U$ on which $F_0$ is
locally constant and $f$ is étale. Write
$j':V_Y=f^{-1}(V)\hookrightarrow Y$. A surjection from a finite sum
of regular permutation modules on the generic fibre spreads, after
shrinking $V$ if necessary, to
\[
 (f|_{V_Y})_*\F_p^r\twoheadrightarrow F_0|_V.
\]
Extension by zero and finite pushforward give a morphism
\[
 P=f_*j'_!\F_p^r=j_!(f|_{V_Y})_*\F_p^r\longrightarrow F_0
\]
whose cokernel $Q$ is supported on finitely many closed points.
The equality of the two sheaves is checked on geometric stalks,
using finiteness of $f$; it is also proper base change for the open
extension-by-zero square. Let $I$ be its image and $N$ its kernel.
The preliminary dimension bound gives $H^3(U,N)=0$, and
$H^q(U,Q)=0$ for $q>0$. The two short exact sequences for $I$
therefore make $H^2(U,F_0)$ a quotient of $H^2(U,P)$.

On $Y$, the complement of $V_Y$ is finite and has residue fields
of $p$-cohomological dimension zero. The localization sequence
therefore gives
\[
 H^2(U,P)=H^2(Y,j'_!\F_p)^r
          \simeq H^2(Y,\F_p)^r=0.
\]
The last equality is Lemma~\ref{A:cyclotomic-Pic}, (iii), and
$\mu_p\simeq\F_p$ on $Y$. Hence $H^2(U,F_0)=0$ for every
constructible $F_0$. The preliminary bound, filtered-colimit
continuity and coefficient d\'evissage give $\cd_p(U)\leq1$ for
arbitrary $p$-primary torsion sheaves.

There are infinitely many closed points on $U$. Choose two,
write $Z$ for their union and $j_0:U-Z\hookrightarrow U$.
Since $U$ is connected, the cokernel of
\[
 H^0(U,\F_p)=\F_p\longrightarrow H^0(Z,\F_p)=\F_p^2
\]
is nonzero and injects into $H^1(U,j_{0!}\F_p)$. Thus $\cd_p(U)=1$,
proving (ii); (ii) implies (i) immediately.

If (iii) fails, the finite-pushforward obstruction applies to every
nonempty affine $U$, whereas the preliminary bound remains two.
Finally $\cd_p(K)=1$, so the required eventual equality of A.2
is precisely equivalent to (i).
\end{proof}

\begin{corollary}[Proved; an exact Picard criterion for the one-dimensional gate]
\label{A:Dedekind-Pic-criterion}
Let $R$ be an excellent Dedekind domain, let $p$ be invertible in $R$,
and suppose $\cd_p(\Frac R)=1$. For any nonempty affine open
$U\subset\Spec R$, the following are equivalent:
\begin{enumerate}[label=(\roman*)]
\item $\cd_p(U)\leq1$;
\item $\operatorname{Pic}(Y)/p=0$ for the normalization $Y$ of $U$
in every finite extension of $\Frac R$.
\end{enumerate}
Consequently A.2 for this $R$ is equivalent to existence of one
$0\ne a\in R$ such that
\[
 \operatorname{Pic}(R_L[1/a])/p=0
 \quad\hbox{for every finite }L/\Frac R,
\]
where $R_L$ is the integral closure of $R$ in $L$.
\end{corollary}
\begin{proof}
Finite extensions have $p$-cohomological dimension at most one.
The divisor and Kummer argument in Lemma~\ref{A:cyclotomic-Pic}
therefore applies to every $Y$ here, giving
$H^2(Y,\mu_p)=\operatorname{Pic}(Y)/p$.
Normalization is finite by excellence, so exact finite pushforward
proves (i)$\Rightarrow$(ii).

For the converse, there are no orderable residue fields when $p=2$:
the regular-local case of Theorem~\ref{A:realexception}, applied to
the discrete valuation ring at such a point, would make its fraction
field orderable and its $2$-cohomological dimension infinite.
The valuation lower bound now gives $\cd_p\kappa(x)=0$ at every
closed point of $U$ and of $Y$. Corollary~\ref{A:generic-plus-dimension}
gives $\cd_p(U)\leq2$. In the permutation-sheaf resolution in the
proof of Theorem~\ref{A:Iwasawa-equivalence}, enlarge $L$ also to
contain $\mu_p$. All its comparisons now apply verbatim: finite
normalization, vanishing of positive boundary cohomology, and
vanishing of third cohomology of the kernel. Hypothesis (ii) gives
$H^2(Y,\F_p)=H^2(Y,\mu_p)=0$, which forces $H^2(U,F_0)=0$
for every constructible $\F_p$-sheaf $F_0$. Continuity and
d\'evissage prove (i).

Every nonempty open of an affine integral scheme contains a nonempty
principal open $D(a)$. Once (ii) holds on $D(a)$, it holds on every
smaller open: normalization commutes with restriction, and the
Picard group of an open Dedekind subscheme is a quotient by the
classes of its omitted primes. Hence all smaller affine opens have
dimension at most one. Proposition~\ref{A:spread} supplies the
opposite bound after one further shrink, simultaneously on all
smaller nonempty affine opens. This is exactly A.2. Its converse
implies (i) on a principal subopen and then (ii), proving the final
equivalence.
\end{proof}

\begin{corollary}[Proved; a single arithmetic instance contains the full odd-prime mu-conjecture]
\label{A:canonical-Iwasawa}
For an odd prime $p$, put
\[
 R_p=\Z[\mu_{p^\infty},1/p],\qquad K_p=\Q(\mu_{p^\infty}).
\]
A.2 for $\Spec R_p$ at $p$ is equivalent to Iwasawa's assertion
$\mu(F_{\mathrm{cyc}}/F)=0$ for every number field $F$.
\end{corollary}
\begin{proof}
The forward implication of Theorem~\ref{A:Iwasawa-equivalence} gives
$\mu(L)=0$ for every finite $L/K_p$. For a number field $F$, take
$L=FK_p$. This is a finite extension of $K_p$, and it is a finite
extension of $F_{\mathrm{cyc}}$ of degree dividing $p-1$.
Extension and norm on the direct limits of $p$-class groups have
composite multiplication by this degree, which is an automorphism
on a $p$-primary torsion group. The $p$-class group of
$F_{\mathrm{cyc}}$ is consequently a direct summand of that of $L$.
The latter is $p$-divisible because $\mu(L)=0$, so the former is
$p$-divisible as well. The structure argument in
Lemma~\ref{A:cyclotomic-Pic} then gives $\mu(F_{\mathrm{cyc}}/F)=0$.
Conversely, every finite $L/K_p$ is the cyclotomic $\Z_p$-extension
of a number field after choosing a finite initial layer. The asserted
vanishing for all number fields gives (iii) and hence A.2.
\end{proof}

\begin{remark}[Literature check and scope]
We searched the exact combinations ``cohomological dimension'',
``cyclotomic'', ``Dedekind'', and ``mu conjecture'', and opened
the primary sources above. The related Galois-group criterion is
\citet[Theorem 11.3.7]{A_NSW}; it concerns the maximal pro-$p$
extension with restricted ramification. Our proof spells out the
additional comparisons for arbitrary constructible sheaves on
$\Spec R$ and for finite extensions ramified at points of that
scheme. We claim no priority. The sources checked in this run do
not prove the general $\mu$-conjecture; Sharifi's Conjecture 3.2.17
states it explicitly. The 2026 abelian and elliptic-curve results
found in the search do not supply its assertion for all number fields.

This leaves that particular cyclotomic instance undecided. The
unrestricted assertion A.2 is refuted independently by the non-Galois
tower in Theorem~\ref{A:A2-persistent-arithmetic-counterexample}; no
claim about the cyclotomic $\mu$-conjecture is inferred from it.
\end{remark}
